\chapter{答案-代数式}

\section{整体带入思维}

\subsection{代整体带入-题目 1}

\label{代整体带入-题目 1} 

暂无答案。

\section{代数最值问题}

\subsection{代数最值问题-题目 1}

\label{代数最值问题-题目 1}  % 为答案设置标签，记录当前页码

\subsection*{已知条件}
\(A, B, C\) 三种原料每袋重量（kg）依次为 \(1, 2, 3\)，每袋价格（万元）依次为 \(3, 2, 5\)。  
生产产品需要 \(A, B, C\) 的袋数分别为 \(x_1, x_2, x_3\)（均为正整数）。  
总重量 \(W = 1\cdot x_1 + 2\cdot x_2 + 3\cdot x_3\)（kg），  
总成本 \(C = 3x_1 + 2x_2 + 5x_3\)（万元）。  
要求 \(W \geq 13\)，求使成本最低的 \((x_1, x_2, x_3)\)。

\subsection*{第一步：写出总重量表达式}
\[
W = x_1 + 2x_2 + 3x_3
\]
故第一空填 \(\boxed{x_1 + 2x_2 + 3x_3}\)。

\subsection*{第二步：分析性价比}
每千克原料的成本：
\[
\frac{3}{1}=3,\quad \frac{2}{2}=1,\quad \frac{5}{3}\approx 1.6667
\]
因此 \(B\) 最便宜，其次 \(C\)，\(A\) 最贵。为降低成本，应尽可能多用 \(B\)，少用 \(A\)。

\subsection*{第三步：枚举可行的正整数组合（从大 \(x_2\) 开始）}
设 \(x_2 = t\)，则已得重量 \(2t\)，还需至少 \(13 - 2t\) 千克（若为负数则已满足）。  
由于 \(x_1, x_3 \geq 1\)，我们需考虑正整数约束。

\subsubsection*{情况 \(t = 6\)}
\(2t = 12\)，还需至少 \(1\) kg。但 \(x_3 \geq 1\) 增加 \(3\) kg，总重至少 \(12+3=15\)，此时 \(x_1\) 至少 \(1\)，取 \(x_1=1, x_3=1\)，得 \(W=1+12+3=16\)，成本 \(C=3+12+5=20\)。

\subsubsection*{情况 \(t = 5\)}
\(2t = 10\)，还需至少 \(3\) kg。取 \(x_3=1\)（加 \(3\) kg）得 \(W=10+3=13\)，但此时 \(x_1=0\) 不满足正整数。故必须增加 \(x_1\)，取 \(x_1=1\)，则 \(W=1+10+3=14\)，成本 \(C=3+10+5=18\)。  
尝试 \(x_3=2\)（加 \(6\) kg）得 \(W=1+10+6=17\)，成本 \(3+10+10=23\)，更高。  
因此 \(t=5\) 时最优为 \((x_1,x_2,x_3)=(1,5,1)\)，成本 \(18\)。

\subsubsection*{情况 \(t = 4\)}
\(2t = 8\)，还需至少 \(5\) kg。组合 \(C\) 和 \(A\)：  
- \(x_3=1\)（3 kg），还需 \(2\) kg，需 \(x_1=2\)，总重 \(2+8+3=13\)，成本 \(6+8+5=19\)。  
- \(x_3=2\)（6 kg），总重 \(8+6=14\)，需 \(x_1=1\)，成本 \(3+8+10=21\)，更高。  
最小成本 \(19\)。

\subsubsection*{情况 \(t = 3\)}
\(2t = 6\)，还需 \(7\) kg。  
- \(x_3=2\)（6 kg），还需 \(1\) kg，\(x_1=1\)，总重 \(1+6+6=13\)，成本 \(3+6+10=19\)。  
- \(x_3=1\)（3 kg），还需 \(4\) kg，\(x_1=4\)，总重 \(4+6+3=13\)，成本 \(12+6+5=23\)。  
最小成本 \(19\)。

\subsubsection*{情况 \(t = 2\)}
\(2t = 4\)，还需 \(9\) kg。  
- \(x_3=3\)（9 kg），总重 \(4+9=13\)，但 \(x_1=0\) 不合法，加 \(x_1=1\) 得 \(W=1+4+9=14\)，成本 \(3+4+15=22\)。  
- \(x_3=2\)（6 kg），还需 \(3\) kg，\(x_1=3\)，总重 \(3+4+6=13\)，成本 \(9+4+10=23\)。  
最小成本 \(22\)。

\subsubsection*{情况 \(t = 1\)}
\(2t = 2\)，还需 \(11\) kg。  
- \(x_3=3\)（9 kg），还需 \(2\) kg，\(x_1=2\)，总重 \(2+2+9=13\)，成本 \(6+2+15=23\)。  
- \(x_3=4\)（12 kg），则总重 \(2+12=14\)，需 \(x_1=1\)，成本 \(3+2+20=25\)。  
最小成本 \(23\)。

\subsection*{第四步：比较所有成本}
\begin{center}
\begin{tabular}{c|c|c|c}
\(x_2\) & \((x_1,x_2,x_3)\) & \(W\) & 成本 \\
\hline
6 & (1,6,1) & 16 & 20 \\
5 & (1,5,1) & 14 & 18 \\
4 & (2,4,1) & 13 & 19 \\
3 & (1,3,2) & 13 & 19 \\
2 & (1,2,3) & 14 & 22 \\
1 & (2,1,3) & 13 & 23 \\
\end{tabular}
\end{center}
可见最小成本为 \(18\)，对应 \(x_1=1,\ x_2=5,\ x_3=1\)。

\subsection*{第五步：填写答案}
总重量表达式：\(x_1 + 2x_2 + 3x_3\)；  
成本最低时的袋数依次为 \(1, 5, 1\)。

\[
\boxed{x_1+2x_2+3x_3}\quad\boxed{1,5,1}
\]


\subsection{代数最值问题-题目 2}

\label{代数最值问题-题目 2}  % 为答案设置标签，记录当前页码


(1) 根据新运算的定义及已知条件：
\[
\begin{cases}
K(1,2)=a+2b=7,\\[2mm]
K(-2,3)=-2a+3b=0.
\end{cases}
\]
解方程组：由第二个方程得 \(3b=2a\)，即 \(a=\dfrac{3}{2}b\)，代入第一个方程：
\[
\frac{3}{2}b+2b=7 \;\Longrightarrow\; \frac{7}{2}b=7 \;\Longrightarrow\; b=2,
\]
于是 \(a=3\)。故 \(a=3,\; b=2\)。

(2) 已知 \(x,y\geq 0\)，且 \(x+2y=10\)，则 \(x=10-2y\)。
由 \(x\geq 0\) 得 \(10-2y\geq 0\)，即 \(y\leq 5\)，结合 \(y\geq 0\) 得 \(0\leq y\leq 5\)。
于是
\[
4x-y = 4(10-2y)-y = 40-8y-y = 40-9y.
\]
因为 \(y\in[0,5]\)，函数 \(40-9y\) 单调递减，所以当 \(y=0\) 时取最大值 \(40\)，当 \(y=5\) 时取最小值 \(40-45=-5\)。故
\[
-5 \leq 4x-y \leq 40.
\]

(3) 由 (1) 知 \(K(y,z)=3y+2z\)，\(K\!\left(x,\dfrac{1}{2}y\right)=3x+2\cdot\dfrac{1}{2}y=3x+y\)。
根据已知条件：
\[
\begin{cases}
3y+2z = 3+x,\\[2mm]
3x+y = 4-3x.
\end{cases}
\]
由第二个方程得 \(y=4-6x\)；代入第一个方程：
\[
3(4-6x)+2z = 3+x \;\Longrightarrow\; 12-18x+2z = 3+x \;\Longrightarrow\; 2z = 9+19x \;\Longrightarrow\; z = \frac{19x-9}{2}.
\]
因为 \(x,y,z\) 均为非负数，所以
\[
\begin{cases}
x \geq 0,\\
y = 4-6x \geq 0 \;\Longrightarrow\; x \leq \dfrac{2}{3},\\
z = \dfrac{19x-9}{2} \geq 0 \;\Longrightarrow\; x \geq \dfrac{9}{19}.
\end{cases}
\]
因此 \(x\) 的取值范围是 \(\dfrac{9}{19}\leq x \leq \dfrac{2}{3}\)。
所求表达式为
\[
x-3y+4z = x-3(4-6x)+4\cdot\frac{19x-9}{2}
= x-12+18x+2(19x-9)
= 19x-12+38x-18 = 57x-30.
\]
该函数随 \(x\) 增大而增大，故
\[
\text{最小值}=57\cdot\frac{9}{19}-30 = 3\cdot9-30 = -3,\qquad
\text{最大值}=57\cdot\frac{2}{3}-30 = 38-30 = 8.
\]
所以 \(x-3y+4z\) 的取值范围是 \([-3,\,8]\)。


\subsection{代数最值问题-题目3-2023七下·石景山期末}

\label{代数最值问题-题目3-2023七下·石景山期末}  

已知方程 \(x+y=3\)，对于它的任意一个解 \((x,y)\)，定义“关联值”为：
\[
\text{关联值}=
\begin{cases}
x-y, & x\ge y,\\
y-x, & x<y,
\end{cases}
\]
即关联值就是 \(|x-y|\).

(1) 写出方程 \(x+y=3\) 的一个解，并指明此时方程的“关联值”

取 \(x=1,\;y=2\)，则 \(x+y=3\) 成立，且 \(x<y\)，所以关联值为 \(y-x=2-1=1\).
（答案不唯一，如 \(x=2,y=1\) 关联值也为 1；\(x=0,y=3\) 关联值为 3 等）

(2) 若“关联值”为 4，写出所有满足条件的解

关联值为 4 即 \(|x-y|=4\)，联立 \(x+y=3\)：
\[
\begin{cases}
x+y=3,\\
|x-y|=4.
\end{cases}
\]
分两种情况：
\begin{itemize}
\item 当 \(x-y=4\) 时，解得 \(x=\dfrac{7}{2},\;y=-\dfrac{1}{2}\)；
\item 当 \(x-y=-4\) 时，解得 \(x=-\dfrac{1}{2},\;y=\dfrac{7}{2}\).
\end{itemize}
故满足条件的解为 \(\left(\dfrac{7}{2},-\dfrac{1}{2}\right)\) 和 \(\left(-\dfrac{1}{2},\dfrac{7}{2}\right)\).

(3) 直接写出最小“关联值”及当关联值为 1 时 \(x\) 的取值范围

由 \(x+y=3\) 得 \(|x-y|=|2x-3|\)（或 \(|3-2y|\)）。  
当 \(x=y=\dfrac{3}{2}\) 时，关联值为 0，且显然关联值非负，故最小关联值为 \(\boxed{0}\).

当关联值为 1 时，\(|x-y|=1\)，结合 \(x+y=3\) 得：
\[
\begin{cases}
x+y=3,\\
x-y=1
\end{cases}
\quad\text{或}\quad
\begin{cases}
x+y=3,\\
x-y=-1.
\end{cases}
\]
解得 \((x,y)=(2,1)\) 或 \((1,2)\).  
因此 \(x\) 只能取 \(1\) 或 \(2\)，即 \(x\) 的取值范围是 \(\{1,2\}\)（或写成 \(x=1\) 或 \(x=2\)）.

\vspace{2mm}
\noindent\textbf{答案：}
\begin{enumerate}
\item 解：\((x,y)=(1,2)\)，关联值为 1；（答案不唯一）
\item \(\left(\dfrac{7}{2},-\dfrac{1}{2}\right),\;\left(-\dfrac{1}{2},\dfrac{7}{2}\right)\)；
\item 最小关联值为 0；当关联值为 1 时，\(x=1\) 或 \(x=2\).
\end{enumerate}

\subsection{代数最值问题-题目3-2025北京三帆中学初一下期中第2题}

\label{代数最值问题-题目3-2025北京三帆中学初一下期中第2题}  


已知方程组
\[
\begin{cases}
x + 2y = m + 5 \quad &(1)\\
2x + y = m + 4 \quad &(2)
\end{cases}
\]
其中 \(x, y \geq 0\)，\(S = x - y\)。

\textbf{第一步：解方程组，用 \(m\) 表示 \(x, y\)。}

将 \((2) \times 2\) 减去 \((1)\)：
\[
(4x + 2y) - (x + 2y) = 2(m+4) - (m+5)
\]
\[
3x = 2m + 8 - m - 5 = m + 3 \quad\Rightarrow\quad x = \frac{m+3}{3}.
\]

代入 \((1)\)：
\[
\frac{m+3}{3} + 2y = m + 5 \quad\Rightarrow\quad 2y = m + 5 - \frac{m+3}{3} = \frac{3m+15 - m - 3}{3} = \frac{2m+12}{3}
\]
\[
y = \frac{m+6}{3}.
\]

\textbf{第二步：利用非负条件求 \(m\) 的范围。}
\[
x \geq 0 \;\Rightarrow\; \frac{m+3}{3} \geq 0 \;\Rightarrow\; m \geq -3,
\]
\[
y \geq 0 \;\Rightarrow\; \frac{m+6}{3} \geq 0 \;\Rightarrow\; m \geq -6.
\]
取交集得 \(m \geq -3\)。

\textbf{第三步：计算 \(S = x - y\)。}
\[
S = \frac{m+3}{3} - \frac{m+6}{3} = \frac{(m+3)-(m+6)}{3} = \frac{-3}{3} = -1.
\]
\(S\) 为常数，与 \(m\) 无关。因此 \(S\) 的取值范围是 \(\boxed{-1}\)。


